\section{Conclusion}
\label{sec:conclusion}

This dissertation asked a single system-level question: when motion prediction
and risk-aware planning are coupled, which apparent component improvements
remain visible in the executed give-way behaviour? The answer is more
conditional than an isolated prediction or planning benchmark would suggest.
Simple task adaptation produces a large and directionally consistent
in-distribution prediction gain, but neither added Transformer complexity nor
adaptive risk establishes uniform system-level superiority. The main
contribution is therefore an evidence chain that follows a model change from
the frozen offline test through calibration, risk allocation, optimisation,
supervision and executed geometry.

\paragraph{H1: supported within the frozen distribution.}
Retraining the MultiPath prediction head reduces rollout-macro NLL from 2.171
to 1.857~nats/step, ADE from 1.283 to 0.100~m and FDE from 2.644 to 0.121~m
relative to the pretrained control. The direction is preserved across all five
held-out initialisation groups and against three physical baselines. This
supports task adaptation for the controlled Town05 give-way distribution, not
generalisation to other maps, interactions or real vehicles.

\paragraph{H2: not supported for the tested Transformer configurations.}
The mean-only Transformer slightly improves upon its matched MLP adapter, while
the distributional Transformer is slightly worse than its matched MLP adapter;
neither exceeds simple task adaptation. Sequence ablations confirm that the
Transformers use their temporal input, so the result is not simple input
neglect. Because the full models are not parameter matched and many runs reach
the epoch ceiling, the conclusion concerns the tested adapters and budget,
not the general value of Transformer architectures.

\paragraph{H3: not supported as a consistent transfer claim.}
The adapted predictor remains better in every in-loop manipulation check, yet
only 2/8 risk-policy--target-style cells jointly improve completion and
physical separation. Better prediction is therefore real but does not have a
uniform driving effect once it enters the coupled stack. Predictor quality
must be evaluated in the planning contexts through which it can change action.

\paragraph{H4: not supported as a universal dominance claim.}
Adaptive risk dominates a fixed comparator in 3/12 prespecified contrasts and
fails to dominate in the other nine. It should therefore be interpreted as a
context-dependent operating point rather than a universal replacement for
fixed risk. The supervisor, solver path and controller acceptance remain part
of this system-level comparison; mechanism telemetry constrains attribution
but does not justify assigning the result to the supervisor alone.

% Provisional tracked Conclusion insertion. The final-only evidence builder atomically
% replaces this file only after every scientific and provenance gate passes.
\paragraph{Supervisor-authority conclusion pending scientific closure.}
The outcome-specific supervisor-authority conclusion is intentionally unavailable
in this non-submission draft. This placeholder contains no outcome estimate,
effect direction, interval interpretation, ranking, superiority statement or
claim that the authority study has reached scientific closure.


Taken together, the four hypotheses refine the project's original expectation
that adaptive risk would simply outperform fixed risk. The stronger conclusion
is not that prediction or adaptive risk is ineffective, but that their benefit
depends on where they alter the closed-loop safety--efficiency frontier. Within
the stated simulator population, the dissertation supplies reproducible
evidence for that coupling and for the importance of retaining negative and
mixed results. It does not claim real-world safety, zero collision risk,
mathematical MPC feasibility or architecture-wide Transformer inferiority.
